\begin{theorem}[Гач, Сипсер]
$\bpp \subset \Sigma_2^p \cap \Pi_2^p$.
\end{theorem}

\begin{proof}
Заметим, что $\bpp$ совпадает с $\mathbf{coBPP}$ в силу симметричности определения сего класса (при переходе к дополнению оно не поменяется). Поэтому достаточно доказать только $\bpp \subset \Sigma_2^p$: из этого будет следовать, что $\mathbf{coBPP} \subset \Pi_2^p$, то есть $\bpp \subset \Pi_2^p$.

Пусть $A \in \bpp$. Рассмотрим алгоритм $V$, дающий для слов длины $n$ правильный ответ с вероятностью $1 - \epsilon$, где $\epsilon < \frac{1}{2^n}$. Пусть он требует $m = \mathrm{poly}(n)$ случайных битов. Для каждого $x$ обозначим через $R_x$ множество таких $r$ (размера $m$), что $V(x, r) = 1$. Если $x \in A$, то $|R_x| > 2^m(1-\epsilon)$, а если $x \not\in A$, то $|R_x| < 2^m\epsilon$. Будем называть сдвигом строки $r$ из $\{0, 1\}^m$ на вектор $s$ побитовое сложение $r \oplus s$ (можно это рассматривать как сдвиг в векторном пространстве размерности $m$ над полем из двух элементов). Также будем говорить о сдвиге множества $R_x$: $R_x \oplus s = \{r \oplus s\colon r \in R_x\}$. Докажем, что для некоторого $k$ при $|R_x| > 2^m(1-\epsilon)$ найдутся сдвиги $s_1, \ldots, s_k$, такие что $R_x \oplus s_i$ в объединении по всем $i$ дают всё $\{0, 1\}^m$, а при $|R_x| < 2^m\epsilon$ таких сдвигов не найдётся.

Для большого $R_x$ воспользуемся простейшим вероятностным методом. Пусть $s_1, \ldots, s_k$ выбраны случайно и независимо. Если сдвиги $R_x$ не покрывают всё пространство, то найдётся $y$, не принадлежащий ни одному из множеств $R_x \oplus s_i$. Это эквивалентно тому, что $R_x$ не содержит ни один из векторов $y \oplus s_1, \ldots, y \oplus s_k$. Эти вектора тоже распределены равномерно и независимо, поэтому для фиксированного $y$ вероятность того, что все они оказались вне $R_x$, равна $\epsilon^k$. Вероятность же того, что такое случится хотя бы для одного $y$, не превосходит $2^m \epsilon^k$. Если $k > \frac{m}{n}$, то эта вероятность будет меньше 1, то есть с положительной вероятностью такого $y$ не найдётся, значит, для некоторого набора сдвигов образы $R_x$ покроют всё пространство, что и требовалось.

Если же размер $R_x$ меньше $2^m\epsilon$, то $k$ его сдвигов покроют меньше $2^m\epsilon k$ элементов, что меньше $2^m$ при $k < 2^n$. Таким образом, при $k = m$ всё пространство быть покрытым не может.

Осталось показать, как указанное свойство можно задать $\Sigma_2$-формулой. Это делается просто:
\[
\exists s_1,\ldots, s_k \forall y\colon \left[V(x, y \oplus s_1) \vee \ldots \vee V(x, y \oplus s_k) \right].
\]
Поскольку $k$ полиномиально, а $V$ работает за полиномиальное время, выражение в скобках также вычисляется за полиномиальное время. Таким образом, $\bpp \subset \Sigma_2^p$, что и требовалось.
\end{proof}